\documentclass[10pt,a4paper]{amsart}
\usepackage{amsmath,amssymb,amsthm,graphicx}
% Retained preamble for the Stacks Project chapter snapshot in this directory.
% The chapter's own preamble.tex is a copy of the shared Stacks Project preamble; the block
% below is sliced MECHANICALLY, line for line, from it. Nothing was added, removed or
% redefined; the wrapper class, the xr-hyper/hyperref wiring and the externaldocument list
% are omitted because the assembler supplies the document class.
%
% lines taken from preamble.tex: 3, 5, 9, 11, 12, 156, 157, 158, ... (51 lines total)
%
% package coverage restored from the chapter's own preamble (2026-10-13): xy, xy-2cell

\usepackage[all]{xy}
\xyoption{2cell}
\UseAllTwocells

\newenvironment{reference}{\comment}{\endcomment}
\newenvironment{slogan}{\comment}{\endcomment}
\newenvironment{history}{\comment}{\endcomment}
\theoremstyle{plain}
\newtheorem{theorem}[subsection]{Theorem}
\newtheorem{proposition}[subsection]{Proposition}
\newtheorem{lemma}[subsection]{Lemma}
\theoremstyle{definition}
\newtheorem{definition}[subsection]{Definition}
\newtheorem{example}[subsection]{Example}
\newtheorem{exercise}[subsection]{Exercise}
\newtheorem{situation}[subsection]{Situation}
\theoremstyle{remark}
\newtheorem{remark}[subsection]{Remark}
\newtheorem{remarks}[subsection]{Remarks}
\numberwithin{equation}{subsection}
\def\lim{\mathop{\mathrm{lim}}\nolimits}
\def\colim{\mathop{\mathrm{colim}}\nolimits}
\def\Spec{\mathop{\mathrm{Spec}}}
\def\Hom{\mathop{\mathrm{Hom}}\nolimits}
\def\Ext{\mathop{\mathrm{Ext}}\nolimits}
\def\SheafHom{\mathop{\mathcal{H}\!\mathit{om}}\nolimits}
\def\SheafExt{\mathop{\mathcal{E}\!\mathit{xt}}\nolimits}
\def\Sch{\mathit{Sch}}
\def\Mor{\mathop{\mathrm{Mor}}\nolimits}
\def\Ob{\mathop{\mathrm{Ob}}\nolimits}
\def\Sh{\mathop{\mathit{Sh}}\nolimits}
\def\NL{\mathop{N\!L}\nolimits}
\def\CH{\mathop{\mathrm{CH}}\nolimits}
\def\proetale{{pro\text{-}\acute{e}tale}}
\def\etale{{\acute{e}tale}}
\def\QCoh{\mathit{QCoh}}
\def\Ker{\mathop{\mathrm{Ker}}}
\def\Im{\mathop{\mathrm{Im}}}
\def\Coker{\mathop{\mathrm{Coker}}}
\def\Coim{\mathop{\mathrm{Coim}}}
\DeclareMathSymbol{\boxtimes}{\mathbin}{AMSa}{"02}
\def\QCohstack{\mathcal{QC}\!\mathit{oh}}
\def\Cohstack{\mathcal{C}\!\mathit{oh}}
\def\Spacesstack{\mathcal{S}\!\mathit{paces}}
\def\Quotfunctor{\mathrm{Quot}}
\def\Hilbfunctor{\mathrm{Hilb}}
\def\Curvesstack{\mathcal{C}\!\mathit{urves}}
\def\Polarizedstack{\mathcal{P}\!\mathit{olarized}}
\def\Complexesstack{\mathcal{C}\!\mathit{omplexes}}
\def\Pic{\mathop{\mathrm{Pic}}\nolimits}
\def\Picardstack{\mathcal{P}\!\mathit{ic}}
\def\Picardfunctor{\mathrm{Pic}}
\def\Deformationcategory{\mathcal{D}\!\mathit{ef}}
\usepackage{cleveref}
\crefname{originallemma}{original lemma}{original lemmas}
\Crefname{originallemma}{Original lemma}{Original lemmas}
\makeatletter
\newlabel{divisors-lemma-characterize-OD}{{01X0}{source}{}{sourcepointer-1}{}}
\newlabel{divisors-lemma-characterize-OD@cref}{{[originallemma][0][]01X0}{[1][0][]source}{}{}{}}
\makeatother
\crefname{originallemma}{original lemma}{original lemmas}
\Crefname{originallemma}{Original lemma}{Original lemmas}
\makeatletter
\newlabel{lemma-different-ramification}{{0BW9}{source}{}{sourcepointer-2}{}}
\newlabel{lemma-different-ramification@cref}{{[originallemma][0][]0BW9}{[1][0][]source}{}{}{}}
\makeatother
\crefname{originallemma}{original lemma}{original lemmas}
\Crefname{originallemma}{Original lemma}{Original lemmas}
\makeatletter
\newlabel{lemma-norm-different-in-discriminant}{{0BW8}{source}{}{sourcepointer-3}{}}
\newlabel{lemma-norm-different-in-discriminant@cref}{{[originallemma][0][]0BW8}{[1][0][]source}{}{}{}}
\makeatother
\title{Stacks, Discriminants: lemma-norm-different-is-discriminant}
\author{The Stacks Project Authors}
\date{stacks-project:discriminant.tex:local-snapshot-20261012; GNU FDL 1.2 or later}
\begin{document}
\maketitle
\noindent\textit{Editorial note.} Statement span is lines 1941-1955 of sources/stacks-discriminant/pinned-original-source.tex; proof span is lines 1957-1989, the complete author proof (31 body lines).\par\medskip
\noindent\textit{Editorial note.} Navigation stubs: 3 cited-but-not-included labels are supplied as non-load-bearing pointers with their pinned tags where the snapshot tag index has them.\par\medskip
\noindent\textit{Editorial note.} extra\_packages is empty on purpose: the retained preamble already loads xy with options and a second declaration clashes in the assembler wrapper.\par\medskip
\noindent\textit{Editorial note.} Text is retained verbatim, including the author's own typos and placeholders.\par\medskip
\noindent\textit{Editorial note.} Queue index 14 of queue\_discriminant.json (pb=1502, stmt 1941-1955, proof 1957-1989).\par\medskip
\noindent\textit{Allowed standard prerequisites (editorial).}\begin{itemize}
\item Explicit prerequisite boundary. The cited inputs are prerequisites; the argument in the retained span is not replaced by any cited result.
\end{itemize}\medskip
\ifdefined\hypertarget\hypertarget{sourcepointer-1}{}\fi
\noindent\textit{Source reference.} Original lemma 01X0: sources/stacks-discriminant/pinned-original-source.tex, source lines 1989--1989. Cited by the retained statement or proof but not part of the collected spans; supplied only as a navigation pointer. The argument does not depend on its content.\par\medskip
\ifdefined\hypertarget\hypertarget{sourcepointer-2}{}\fi
\noindent\textit{Source reference.} Original lemma 0BW9: sources/stacks-discriminant/pinned-original-source.tex, source lines 1907--1912. Cited by the retained statement or proof but not part of the collected spans; supplied only as a navigation pointer. The argument does not depend on its content.\par\medskip
\ifdefined\hypertarget\hypertarget{sourcepointer-3}{}\fi
\noindent\textit{Source reference.} Original lemma 0BW8: sources/stacks-discriminant/pinned-original-source.tex, source lines 1863--1867. Cited by the retained statement or proof but not part of the collected spans; supplied only as a navigation pointer. The argument does not depend on its content.\par\medskip
\begin{lemma}
\label{lemma-norm-different-is-discriminant}
Let $f : Y \to X$ be a flat quasi-finite morphism of Noetherian schemes.
Let $R \subset Y$ be the closed subscheme defined by $\mathfrak{D}_f$.
\begin{enumerate}
\item If $\omega_{Y/X}$ is invertible,
then $R$ is a locally principal closed subscheme of $Y$.
\item If $\omega_{Y/X}$ is invertible and $f$ is finite, then
the norm of $R$ is the discriminant $D_f$ of $f$.
\item If $\omega_{Y/X}$ is invertible and $f$
is \'etale at the associated points of $Y$, then $R$
is an effective Cartier divisor and there is an
isomorphism $\mathcal{O}_Y(R) = \omega_{Y/X}$.
\end{enumerate}
\end{lemma}

\begin{proof}
Proof of (1). We may work locally on $Y$, hence we may assume
$\omega_{Y/X}$ is free of rank $1$. Say $\omega_{Y/X} = \mathcal{O}_Y\lambda$.
Then we can write $\tau_{Y/X} = h \lambda$ and then we see that
$R$ is defined by $h$, i.e., $R$ is locally principal.

\medskip\noindent
Proof of (2). We may assume $Y \to X$ is given by a finite free ring
map $A \to B$ and that $\omega_{B/A}$ is free of rank $1$ as $B$-module.
Choose a $B$-basis element $\lambda$ for $\omega_{B/A}$ and write
$\text{Trace}_{B/A} = b \cdot \lambda$ for some $b \in B$.
Then $\mathfrak{D}_{B/A} = (b)$ and $D_f$ is cut out by
$\det(\text{Trace}_{B/A}(b_ib_j))$ where $b_1, \ldots, b_n$ is a
basis of $B$ as an $A$-module. Let $b_1^\vee, \ldots, b_n^\vee$
be the dual basis.
Writing $b_i^\vee = c_i \cdot \lambda$ we see that
$c_1, \ldots, c_n$ is a basis of $B$ as well.
Hence with $c_i = \sum a_{il}b_l$ we see that $\det(a_{il})$
is a unit in $A$. Clearly,
$b \cdot b_i^\vee = c_i \cdot \text{Trace}_{B/A}$
hence we conclude from the computation in the proof of
Lemma \ref{lemma-norm-different-in-discriminant}
that $\text{Norm}_{B/A}(b)$ is a unit times
$\det(\text{Trace}_{B/A}(b_ib_j))$.

\medskip\noindent
Proof of (3). In the notation above we see from
Lemma \ref{lemma-different-ramification} and the assumption
that $h$ does not vanish in
the associated points of $Y$, which implies that $h$ is a nonzerodivisor.
The canonical isomorphism sends $1$ to $\tau_{Y/X}$, see
Divisors, Lemma \ref{divisors-lemma-characterize-OD}.
\end{proof}

\end{document}
